\documentclass[11pt,a4paper]{article}
\usepackage{amsmath,amssymb,graphicx,geometry,hyperref,fontspec}
\geometry{margin=2.2cm}
\setlength{\parindent}{0pt}\setlength{\parskip}{4pt}
\title{LieRing Beam Physics Manual\\\large symplectic ring modelling with Lie maps (v0.4)}
\date{\today}
\begin{document}
\maketitle
\tableofcontents
\clearpage
\section{LieRing Beam Physics Manual}

LieRing is a modular Julia program for symplectic modelling of a storage ring
(here: the Nuclotron-based lattice \texttt{magnetic\_2.opt}, deuteron, 270 MeV) built
on Hamiltonians and Lie algebra. Every map in the code is derived from the
element Hamiltonians written in this manual; no external code was used for the
physics. OptiM and COSY INFINITY results are used only as external numeric checks.

This manual describes \textbf{only what exists in the code} (status: v0.4). Each
section names the source file that implements it (see Code map).

Contents:

\begin{enumerate}
\item Coordinates and reference particle
\item Coordinate transformation user <-> canonical
\item Element Hamiltonians and Lie maps
\item Linear optics
\item Second-order dispersion, momentum compaction and slip factor
\item Chromaticity
\item RF cavity of finite length
\item Acceleration and adiabatic damping
\item Longitudinal motion, separatrix and synchrotron tune
\item Symplecticity checks and Taylor-map comparison
\item Sextupoles and the optional optimiser
\item Assumptions and hypotheses
\item Code map
\item Programmer manual (stub)
\end{enumerate}

All scripts are run from \texttt{Julia/ring} as \texttt{julia --project=. scripts/NN\_name.jl};
the numbers quoted below are copied from their printed output.

\clearpage
\section{Coordinates and reference particle}

\emph{Implemented in \texttt{src/Constants.jl}.}

The reference particle has charge $Z e$, mass $m$, momentum $p_0$,
velocity $v_0=\beta_0 c$ and Lorentz factor $\gamma_0$. The independent
variable is the arc length $s$ along the reference orbit.

\textbf{User coordinates} (input/output, plots):

\[
u = (x,\ x',\ y,\ y',\ \phi,\ \delta),\qquad x'=\frac{dx}{ds},\quad \delta=\frac{p-p_0}{p_0}.
\]

$\delta$ is the relative \emph{momentum} deviation (not energy). $\phi$ is the RF
phase of arrival relative to the reference particle with the sign convention

\[
\phi = -h\,\omega_0\,\Delta t,\qquad \Delta t = t - t_0(s),
\]

so $\phi>0$ means the particle arrives \emph{early}. $h$ is the harmonic number and
$\omega_0 = 2\pi v_0/C$ the angular revolution frequency.

\textbf{Canonical coordinates} (all maps, Hamiltonians, symplecticity):

\[
z = (x,\ p_x,\ y,\ p_y,\ l,\ \delta),\qquad p_{x,y}=\frac{P_{x,y}}{p_0},\quad l = -v(\delta)\,\Delta t .
\]

$(l,\delta)$ is a canonical pair: $l$ is the path-length-like lag measured
with the particle's own velocity $v(\delta)$, with

\[
\frac{v(\delta)}{v_0}=\frac{1+\delta}{\sqrt{1+\beta_0^2\left((1+\delta)^2-1\right)}} .
\]

The energy term of every Hamiltonian is

\[
e(\delta)=\frac{1}{\beta_0}\sqrt{(1+\delta)^2+\frac{1}{\beta_0^2\gamma_0^2}}-\frac{1}{\beta_0^2},
\]

(the constant is irrelevant for the dynamics and makes $e(0)$-terms vanish after
removing the constant part). Particle data: deuteron, $Z=1$,
$mc^2 = 2\times 938.272$ MeV (as in OptiM).

\clearpage
\section{Coordinate transformation user <-> canonical}

\emph{Implemented in \texttt{src/Constants.jl} (\texttt{canonical\_to\_user}, \texttt{user\_to\_canonical}),
used in \texttt{src/TaylorMap.jl} (\texttt{user\_map}) and \texttt{src/Acceleration.jl}; checked by
\texttt{scripts/06\_coordinates.jl}.}

\subsection{Exact formulas}

Canonical $\to$ user:

\[
p_s=\sqrt{(1+\delta)^2-p_x^2-p_y^2},\qquad x'=\frac{p_x}{p_s},\quad y'=\frac{p_y}{p_s},\qquad
\phi=\frac{2\pi h}{C}\,\frac{l}{v(\delta)/v_0}.
\]

User $\to$ canonical:

\[
p_x=\frac{x'(1+\delta)}{\sqrt{1+x'^2+y'^2}},\quad p_y=\frac{y'(1+\delta)}{\sqrt{1+x'^2+y'^2}},\qquad
l=\frac{C}{2\pi h}\,\frac{v(\delta)}{v_0}\,\phi .
\]

$x$, $y$ and $\delta$ are common to both sets. These follow from
$x'=P_x/P_s$ and from $\phi=-h\omega_0\Delta t$, $l=-v\Delta t$.

\subsection{Applied as full polynomials, not linearised}

Both functions are generic Julia code. In \texttt{user\_map} they are evaluated on the
\texttt{TaylorN} variables, which produces the complete Taylor expansion of the
square roots and of $1/v(\delta)$; the result is truncated only at the map
order (\texttt{MAP\_ORDER = 3}). The map in user coordinates is

\[
U = T\circ M\circ T^{-1},\qquad T = \text{canonical\_to\_user}.
\]

Tracking (\texttt{accelerate}, \texttt{turn\_snapshots}) uses the same functions on
floating-point numbers, i.e. exactly. The code was checked for hidden
linearisations: none were found. (The linear-optics module reports
$D'=dx/d\delta$-slope in $p_x$; to first order $p_x=x'$, so this is not an
approximation at the order it is used.)

Examples printed by \texttt{06\_coordinates.jl} (order 3):

\[
x' = p_x - p_x\delta + \tfrac12 p_x^3 + \tfrac12 p_x p_y^2 + p_x\delta^2+\dots,\qquad
\phi = 0.267624\,l - 0.204533\,l\delta + 0.228642\,l\delta^2+\dots
\]

\subsection{Symplecticity holds only in canonical variables}

$(x,x',y,y',\phi,\delta)$ is \textbf{not} a canonical set ($x'$ depends on
$\delta$ and $p_y$; $\phi$ is scaled by $1/v(\delta)$). Therefore the
symplectic condition $J^TSJ=S$ is checked only for maps in $z$. The
user-coordinate map is the same physics written in other variables.

\subsection{Numeric check (scripts/06\_coordinates.jl)}

\begin{center}\small\begin{tabular}{p{0.475\linewidth}p{0.475\linewidth}}\hline
\textbf{check} & \textbf{result}\\\hline
polynomial round trip $T(T^{-1}(u))-u$, all orders kept & max coefficient $\le 1.1\times10^{-16}$\\
exact functions round trip at $u=(5\,\mathrm{mm},20\,\mathrm{mrad},5\,\mathrm{mm},20\,\mathrm{mrad},1\,\mathrm{rad},10^{-2})$ & $0$\\
order-3 polynomial $T$ vs exact $T$ at that point & $8.7\times10^{-7}$ (truncation)\\
linearised $T$ vs exact $T$ at that point & $7.6\times10^{-3}$\\
$\max\vert J^TSJ-S\vert $ canonical one-turn map at $10^{-3}$ & $6.8\times10^{-6}$ (order-3 truncation)\\
$\max\vert J^TSJ-S\vert $ user one-turn map at $10^{-3}$ & $0.12$ (non-canonical variables)\\
\hline\end{tabular}\end{center}

Selected one-turn coefficients (270 MeV, RF off):

\begin{center}\small\begin{tabular}{p{0.317\linewidth}p{0.317\linewidth}p{0.317\linewidth}}\hline
\textbf{coefficient (canonical / user)} & \textbf{canonical} & \textbf{user}\\\hline
$(x\vert p_x)$ / $(x\vert x')$ & 4.61348258 & 4.61348258\\
$(x\vert \delta)$ & 2.42044391e-2 & 2.42044391e-2\\
$(l\vert \delta)$ / $(\phi\vert \delta)$ & 62.6600041 & 16.7693216\\
$(l\vert \delta^2)$ / $(\phi\vert \delta^2)$ & -71.5275409 & -31.9585141\\
$(l\vert p_x^2)$ / $(\phi\vert x'^2)$ & -86.2567313 & -23.0843723\\
$(x\vert p_x^2\delta)$ / $(x\vert x'^2\delta)$ & 343.450443 & 343.541752\\
\hline\end{tabular}\end{center}

The last line shows a genuinely nonlinear difference between the systems.

\clearpage
\section{Element Hamiltonians and Lie maps}

\emph{Types: \texttt{src/Elements.jl}; Hamiltonians: \texttt{src/Hamiltonians.jl}; Lie operators:
\texttt{src/LieOperators.jl}; factor-by-factor tracking: \texttt{src/Tracking.jl}.}

\subsection{Hamiltonians}

With $p_s=\sqrt{(1+\delta)^2-p_x^2-p_y^2}$, curvature $h=1/\rho$ and
normalised strengths $k_1=B'/(B\rho)$, $k_2=B''/(B\rho)$, the
$s$-Hamiltonians are (no fringe fields, hard-edge, field constant in $s$):

\begin{center}\small\begin{tabular}{p{0.475\linewidth}p{0.475\linewidth}}\hline
\textbf{element} & \textbf{Hamiltonian $H(z)$}\\\hline
Drift & $e(\delta)-p_s$\\
Quadrupole & $e(\delta)-p_s+\tfrac{k_1}{2}(x^2-y^2)$\\
Sextupole & $e(\delta)-p_s+\tfrac{k_2}{6}(x^3-3xy^2)$\\
Sector bend (with gradient) & $e(\delta)-(1+hx)p_s+hx+\tfrac{h^2x^2}{2}+\tfrac{k_1}{2}(x^2-y^2)$\\
Marker & $0$\\
Thin RF kick (\texttt{RFCavity}) & kick only, see RF cavity of finite length\\
Finite cavity (\texttt{Cavity}) & $e(\delta)-p_s+A\cos\Phi(l,\delta,s)$ (exact integrator only; as drift in the maps)\\
\hline\end{tabular}\end{center}

For the bend, $\partial H/\partial x|_{z=0}=-h+h=0$: the reference arc is a
solution. All Hamiltonians are expanded as \texttt{TaylorN} polynomials to
\texttt{JET\_ORDER = 4} in $z$, and the constant term is removed.

\subsection{Map of an element}

For an $s$-independent $H$ over length $L$ the map is the Lie transformation

\[
\mathcal M = \exp(-L:H:),\qquad :f:g = [f,g]=\sum_{i}\left(\frac{\partial f}{\partial q_i}\frac{\partial g}{\partial p_i}-\frac{\partial f}{\partial p_i}\frac{\partial g}{\partial q_i}\right),
\]

applied to the coordinate functions as the series
$\exp(:f:)z=\sum_n :f:^n z/n!$, truncated at \texttt{MAP\_ORDER = 3}.
The map is then Dragt–Finn factorised,

\[
\mathcal M = R\ \circ\ e^{:f_3:}\ \circ\ e^{:f_4:},
\]

$R$ the linear (symplectic) matrix and $f_3,f_4$ homogeneous cubic/quartic
generators. Tracking (\texttt{Tracking.jl}) applies $R$ exactly and each
$e^{:f_k:}$ by integrating the Hamiltonian flow of $-f_k$ with the implicit
midpoint rule, so the tracking map is symplectic to round-off even though
the generators are truncated.

\clearpage
\section{Linear optics}

\emph{\texttt{src/LinearOptics.jl}, \texttt{src/OptiMParser.jl}, script \texttt{01\_linear\_optics.jl}.}

The linear map of each element is the Hessian of its Hamiltonian:
$R=\exp(L\,S\,\mathcal H)$, $\mathcal H_{ij}=\partial^2H/\partial z_i\partial z_j|_0$.
The one-turn matrix $M$ is the product over the lattice. With
$\cos\mu=\tfrac12\mathrm{Tr}\,M_{x}$, the closed Twiss functions are

\[
\beta=\frac{M_{12}}{\sin\mu},\quad \alpha=\frac{M_{11}-M_{22}}{2\sin\mu},
\]

and the closed dispersion solves $(I-M_4)D=M_{\cdot6}$ (4x4 transverse part,
$D=(D,D',D_y,D_y')$). The momentum compaction is

\[
\alpha_c=\frac{1}{\gamma_0^2}-\frac{(l|\delta)_{\rm closed}}{C},
\]

cross-checked by $\alpha_c=\frac1C\oint hD\,ds$. Twiss and dispersion are
propagated in slices along the ring ("closed" check: after one turn they
reproduce the start to $10^{-13}$).

\begin{center}\small\begin{tabular}{p{0.238\linewidth}p{0.238\linewidth}p{0.238\linewidth}p{0.238\linewidth}}\hline
\textbf{quantity} & \textbf{LieRing} & \textbf{reference} & \textbf{source}\\\hline
C [m] & 117.388296 &  & \\
$\beta_x(0)$ [m] & 19.367522 & 19.3675 & OptiM\\
$\beta_y(0)$ [m] & 4.054470 & 4.05447 & OptiM\\
$\alpha_x(0)$ & -2.5191e-4 & -2.51909e-4 & OptiM\\
$D'(0)$ & -5.0616e-4 & -5.0616e-4 & OptiM\\
$D(0)$ [m] & 0.011094 & 0 & OptiM (see Assumptions)\\
$Q_x$ & 2.461720 & frac 0.446284 & COSY\\
$Q_y$ & 2.461337 & frac 0.476682 & COSY\\
$\alpha_c$ & 0.23047028 & 0.23036260 & COSY\\
\hline\end{tabular}\end{center}

\begin{figure}[htbp]\centering\includegraphics[width=\linewidth]{C:/GIT_REPOS/BELL/Julia/ring/docs/src/assets/magnetic_2_twiss.png}\caption{Closed beta functions, dispersion and magnet layout of magnetic\_2 at 270 MeV (01\_linear\_optics.jl). Panels: beta\_x and beta\_y, horizontal dispersion D\_x, and the magnet layout along s.}\end{figure}

The tunes differ from COSY by 0.015 in both planes, with opposite signs. The
reason is open (fringe-field or element-model treatment in the COSY deck is the
main suspect); the OptiM initial Twiss agree to 5 digits.


\clearpage
\section{Second-order dispersion, momentum compaction and slip factor}

\emph{\texttt{src/SecondOrder.jl}, \texttt{src/TaylorMap.jl}, script \texttt{03\_second\_order.jl}.}

The one-turn Taylor map $M$ (order 3, RF off) gives the closed orbit as a
polynomial in $\delta$:

\[
z_\perp(\delta)=D_1\delta+D_2\delta^2+O(\delta^3),\qquad D_1=(D,D'),\ D_2=(D_2,D_2'),
\]

from the fixed-point condition $M_\perp(z_\perp(\delta),\delta)=z_\perp(\delta)$
order by order: $(I-R)D_1=R_{\cdot6}$, $(I-R)D_2=Q(D_1\delta+\delta e_6)|_{\delta^2}$.
$D_2$ is propagated around the ring by pushing the polynomial orbit through
each element map (\texttt{dispersion\_along}).

The path-length change per turn on the closed orbit is
$\Delta l(\delta)=l_1\delta+l_2\delta^2$ ($l$ component of the map evaluated
on $z_\perp(\delta)$). With the velocity expansion
$\frac{v_0}{v(\delta)}=1-c_1\delta-c_2\delta^2+\dots$ ($c_1=1/\gamma_0^2$):

\[
\alpha_c(\delta)=\alpha_0+\alpha_1\delta,\quad \alpha_0=-\frac{l_1}{C}+c_1,\quad \alpha_1=-\frac{l_2}{C}+c_2,
\]

\[
\eta(\delta)=\eta_0+\eta_1\delta,\quad \eta_0=\alpha_0-\frac{1}{\gamma_0^2},\quad
\eta_1=\alpha_1-\alpha_0c_1-c_2+c_1^2 .
\]

$\eta$ is the relative change of the revolution time, $\Delta T/T=\eta_0\delta+\eta_1\delta^2$,
and enters the phase equation $\Delta\phi/\text{turn}=-2\pi h(\eta_0\delta+\eta_1\delta^2)$.

Results (natural lattice, sextupoles 0), from \texttt{05\_rf\_acceleration.jl}
section 4 and \texttt{03\_second\_order.jl}:

\begin{center}\small\begin{tabular}{p{0.317\linewidth}p{0.317\linewidth}p{0.317\linewidth}}\hline
\textbf{quantity} & \textbf{10 MeV} & \textbf{270 MeV}\\\hline
$\alpha_0$ & 0.23047028 & 0.23047028\\
$\alpha_1$ & 0.33900771 & 0.33900771\\
$\eta_0$ & -0.75895642 & -0.53378406\\
$\eta_1$ & 1.10563172 & 1.01720896\\
\hline\end{tabular}\end{center}

$\alpha_0,\alpha_1$ are purely geometric (magnet strengths scale with $p_0$);
the energy dependence of $\eta$ comes only from $c_1,c_2$.

\begin{figure}[htbp]\centering\includegraphics[width=\linewidth]{C:/GIT_REPOS/BELL/Julia/ring/docs/src/assets/sextupoles_dispersion2.png}\caption{Second-order dispersion D2 along the ring for the natural lattice and for the optional optimiser examples (04\_optional\_sextupole\_optimization.jl).}\end{figure}

Second-order dispersion at $s=0$ (\texttt{03\_second\_order.jl}): $D=0.01109398$ m,
$D'=-5.0616\times10^{-4}$, $D_2=-10.054932$ m, $D_2'=0.459066$; along the ring
$D_2\in[-10.08,\,14.41]$ m; closure $D_2(C)-D_2(0)=-3.7\times10^{-14}$.
Exact-tracking cross-check (closed orbit at fixed $\delta$, cubic fit):
$x_{co}=0.01108483\,\delta-10.054867\,\delta^2$, $\Delta C/C=0.23047050\,\delta+0.339007\,\delta^2$.
Here $\alpha_1$ is defined by $\Delta C/C=\alpha_0\delta+\alpha_1\delta^2$
(the local derivative is $\alpha_0+2\alpha_1\delta$).


\clearpage
\section{Chromaticity}

\emph{\texttt{src/SecondOrder.jl} (\texttt{chromaticity\_orders}), \texttt{src/TaylorMap.jl}
(\texttt{tunes\_at\_delta}), script \texttt{03\_second\_order.jl}.}

For fixed $\delta$ the transverse motion around the closed orbit
$z_\perp(\delta)$ is linearised: the 4x4 Jacobian $J(\delta)$ of the map at
the fixed point gives $Q_{x,y}(\delta)$ from $\cos 2\pi Q=\tfrac12\mathrm{Tr}J$.
Since $\delta$ is a constant of motion without RF, the tunes are
invariants of motion and

\[
Q(\delta)=Q_0+\xi_1\delta+\tfrac12\xi_2\delta^2,\qquad \xi_1=\frac{dQ}{d\delta},\ \xi_2=\frac{d^2Q}{d\delta^2},
\]

obtained by central differences with $\Delta\delta=10^{-4}$ on the polynomial map.
Cross-check: tracking with the exact integrator at several $\delta$, tune
from the least-squares fit $x_{n+1}+x_{n-1}=2\cos(2\pi Q)x_n+b$ (insensitive to a
small fixed-point offset), then a parabola fit of $Q(\delta)$.
Numbers: see the output of \texttt{03\_second\_order.jl} (reproduced in
Sextupoles and the optional optimiser).

\begin{figure}[htbp]\centering\includegraphics[width=\linewidth]{C:/GIT_REPOS/BELL/Julia/ring/docs/src/assets/chromaticity_natural.png}\caption{Tunes versus delta for the natural lattice: map (line) and tracking (markers).}\end{figure}

\begin{center}\small\begin{tabular}{p{0.238\linewidth}p{0.238\linewidth}p{0.238\linewidth}p{0.238\linewidth}}\hline
\textbf{} & \textbf{map} & \textbf{tracking fit} & \textbf{COSY (reference only)}\\\hline
$Q_{x0}$ (frac) & 0.46172014 & 0.46171903 & \\
$\xi_{1x}$ & -1.461538 & -1.461531 & -1.840801\\
$\xi_{2x}$ & 8.5547 & 8.5546 & \\
$Q_{y0}$ (frac) & 0.46133732 & 0.46133566 & \\
$\xi_{1y}$ & -2.225488 & -2.225479 & -1.847679\\
$\xi_{2y}$ & 8.1041 & 8.1038 & \\
\hline\end{tabular}\end{center}

The difference to COSY is open (its chromaticity convention was not verified).

\clearpage
\section{RF cavity of finite length}

\emph{\texttt{src/ExactIntegrator.jl} (\texttt{\_field\_cav}, \texttt{cavity\_parameters},
\texttt{transit\_factor}), type \texttt{Cavity} in \texttt{src/Elements.jl}; checks in
\texttt{05\_rf\_acceleration.jl} section 2.}

The cavity is a \textbf{finite-length} element of length $L$ with a purely
longitudinal on-axis field; it is \emph{not} a drift plus a thin kick. With
$q A_s/p_0=-A\cos(\omega t+\psi)$, $E_s=-\partial A_s/\partial t$ and
$t=s/v_0-l/v(\delta)$, the $s$-Hamiltonian is

\[
H_{\rm cav}=e(\delta)-p_s+A\cos\Phi,\qquad
\Phi=\frac{\omega(s-L/2)}{v_0}-\frac{\omega\,l}{v(\delta)}-\phi_s,
\]

\[
A=\frac{ZV}{p_0c}\,\frac{c}{\omega L},\qquad \omega=h\omega_0=\frac{2\pi h v_0}{C}.
\]

Hamilton's equation gives

\[
\frac{d\delta}{ds}=-\frac{\partial H}{\partial l}=-\frac{ZV}{p_0cL}\frac{c}{v}\sin\Phi,
\]

i.e. $dE/ds=-(ZV/L)\sin\Phi$. Integrated over the cavity for a particle
with $\phi$ the energy gain is

\[
\Delta E=ZVT\sin(\phi_s+\phi),\qquad T=\frac{\sin(\theta/2)}{\theta/2},\ \theta=\frac{\omega L}{v_0}=\frac{2\pi hL}{C}.
\]

$H_{\rm cav}$ depends explicitly on $s$; it is integrated with the implicit
midpoint rule composed to 4th order (Yoshida), which is symplectic for
$s$-dependent Hamiltonians. Transverse RF fields are neglected.

Checks at 270 MeV ($L=1$ m, $V=0.2$ MV, $h=5$, $\phi_s=150^\circ$,
$T=0.99702$):

\begin{center}\small\begin{tabular}{p{0.475\linewidth}p{0.475\linewidth}}\hline
\textbf{check} & \textbf{result}\\\hline
synchronous gain vs $ZVT\sin\phi_s$ & 0.099700856 vs 0.099701840 MeV\\
thin limit $L=0.1/0.01/0.001$ m vs thin kick ($\phi=0.05$) & $-6.3\times10^{-9}$, $-1.5\times10^{-10}$, $-9.9\times10^{-12}$ in $\delta$\\
symplectic error of the cavity Jacobian & $4.4\times10^{-16}$\\
\hline\end{tabular}\end{center}

The thin kick (\texttt{RFCavity}, \texttt{Tracking.rf\_kick}) applies the same energy gain $ZV\sin(\phi_s+\phi)$ at one point; the finite cavity converges to it as $L\to0$ (table above).


\clearpage
\section{Acceleration and adiabatic damping}

\emph{\texttt{src/ExactIntegrator.jl} (\texttt{renormalize}), \texttt{src/Acceleration.jl}
(\texttt{accelerate}, \texttt{turn\_snapshots}), script \texttt{05\_rf\_acceleration.jl}.}

\subsection{Reference momentum update}

All particles, including the synchronous particle (\#1, initially $z=0$),
are tracked through the full lattice with the exact integrator. Right after
the cavity the synchronous particle has $\delta_s\neq0$; it is made the new
reference:

\[
p_0'=p_0(1+\delta_s),\quad 1+\delta'=\frac{1+\delta}{1+\delta_s},\quad
p_x'=\frac{p_x}{1+\delta_s},\ p_y'=\frac{p_y}{1+\delta_s},\quad
\Delta t'=\Delta t-\Delta t_s,\ l'=-v\,\Delta t' .
\]

Positions $x,y$ are unchanged. The magnets are assumed to ramp with $p_0$
(normalised $h,k_1,k_2$ fixed) and the RF frequency follows
$h\,\omega_0(p_0)$ (locked to the synchronous particle).

\subsection{Adiabatic damping}

The physical transverse momentum $P_x=p_0p_x$ is not changed by a
longitudinal field, so $p_x\propto1/p_0$. Between kicks the motion is a
betatron oscillation with fixed optics; the Courant–Snyder action in
normalised variables $J=\tfrac12(\gamma x^2+2\alpha xp_x+\beta p_x^2)$ is an
adiabatic invariant of the \emph{unnormalised} phase space,
$\oint P_x\,dx=p_0\oint p_x\,dx=2\pi p_0J=\text{const}$. Hence
$J\propto1/p_0$ and the amplitude $A=\sqrt{2\beta J}$ scales as

\[
\frac{A(n)}{A(0)}=\sqrt{\frac{p_{\rm inj}}{p_0(n)}} .
\]

Ramp results ($10\to270$ MeV, $V=0.2$ MV, $h=5$, $\phi_s=150^\circ$, then
1000 turns with $\phi_s=180^\circ$): see the table in
Longitudinal motion, separatrix and synchrotron tune.

\clearpage
\section{Longitudinal motion, separatrix and synchrotron tune}

\emph{\texttt{src/SecondOrder.jl} (\texttt{LongitudinalParams}, \texttt{long\_H}, \texttt{fixed\_points},
\texttt{separatrix}, \texttt{bucket\_height}, \texttt{synchrotron\_tune}), plots in
\texttt{src/Plotting.jl}, script \texttt{05\_rf\_acceleration.jl}.}

\subsection{Turn-by-turn map and smooth Hamiltonian}

With one cavity per turn, the motion of $(\phi,\delta)$ per turn $n$ is

\[
\phi_{n+1}=\phi_n-2\pi h\left(\eta_0\delta_{n+1}+\eta_1\delta_{n+1}^2\right),\qquad
\delta_{n+1}=\delta_n+A\left(\sin(\phi_s+\phi_n)-\sin\phi_s\right),\quad A=\frac{ZVT}{\beta_0^2E_0}.
\]

(The tracking itself does not use this map; it integrates the exact
Hamiltonians element by element. The map and the Hamiltonian below are used
for formulas and for drawing the separatrix.) Treating $n$ as continuous,

\[
H(\phi,\delta)=-2\pi h\left(\frac{\eta_0\delta^2}{2}+\frac{\eta_1\delta^3}{3}\right)+A\left(\cos(\phi_s+\phi)+\phi\sin\phi_s\right).
\]

Fixed points on $\delta=0$: stable $\phi=0$, unstable
$\phi_u=\pi-2\phi_s$; $\eta_1$ adds a second pair at
$\delta\approx-\eta_0/\eta_1$ (far outside the bucket here, 0.69 at 10 MeV).
The \textbf{separatrix} is the level line $H(\phi,\delta)=H(\phi_u,0)$; the bucket
is the region around $(0,0)$ inside it. Stability requires
$\eta_0\cos\phi_s>0$: below transition ($\eta_0<0$) and with our sign of
$\phi$ this means $90^\circ<\phi_s\le180^\circ$; acceleration uses
$\phi_s=150^\circ$ (gain $ZVT\sin\phi_s>0$), the flat top $\phi_s=180^\circ$.

\subsection{Synchrotron tune}

Small amplitude:

\[
Q_s^{\rm smooth}=\sqrt{\frac{h\,\eta_0\,ZVT\cos\phi_s}{2\pi\beta_0^2E_0}},\qquad
\cos(2\pi Q_s^{\rm kick})=1-\frac{(2\pi Q_s^{\rm smooth})^2}{2}\ \ \text{(one kick per turn)}.
\]

\begin{center}\small\begin{tabular}{p{0.238\linewidth}p{0.238\linewidth}p{0.238\linewidth}p{0.238\linewidth}}\hline
\textbf{$E$ [MeV], $\phi_s=180^\circ$} & \textbf{tracking} & \textbf{smooth} & \textbf{kick}\\\hline
10 & 0.0784653 & 0.0777018 & 0.0784950\\
270 & 0.0129407 & 0.0129376 & 0.0129411\\
\hline\end{tabular}\end{center}

At 10 MeV $Q_s$ is large and the discrete-kick formula is needed (1\% vs
0.04\%).

\subsection{Ramp 10 -> 270 MeV}

Parameters: $L=1$ m, $V=0.2$ MV, $h=5$ ($f_{RF}$ 1.313 -> 6.200 MHz),
$\phi_s=150^\circ$; gain 0.099702 MeV/turn; 270 MeV reached at turn 2608,
then 1000 turns at $\phi_s=180^\circ$. Bucket half-height at injection
$+0.02425/-0.02369$ (asymmetric because of $\eta_1$), stationary bucket at
270 MeV $0.00976$.

\begin{center}\small\begin{tabular}{p{0.475\linewidth}p{0.475\linewidth}}\hline
\textbf{result} & \textbf{value}\\\hline
tracked & 64 particles x 3608 turns (330 s)\\
mean gain per turn & 0.099698 MeV (expected 0.099702)\\
x amplitude 2 mm -> & 0.8701 mm (ratio 0.43506)\\
y amplitude 2 mm -> & 0.8658 mm (ratio 0.43291)\\
prediction $\sqrt{p_{inj}/p_{fin}}$ & 0.43142\\
ensemble particles lost from the bucket & 0 of 60\\
longitudinal invariant (turn 300 -> 2493) & 1.0000 -> 0.9508\\
\hline\end{tabular}\end{center}

The first \~{}200 turns are not adiabatic (gain = 1\% of $E_{kin}$ per turn at
10 MeV), which explains most of the 5\% drift of the longitudinal invariant.

Ensemble: 60 particles with $\phi\sim N(0,0.6)$ rad,
$\delta\sim N(0,0.25\,\hat\delta_{\rm bucket})$, kept if inside the injection
separatrix, on their dispersion orbit, $x'=y'=0$. This is deliberately
mismatched to the bucket so that filamentation is visible. $\eta_1$ and
$\alpha_1$ enter the separatrix; $\alpha_0=0.23047028$,
$\alpha_1=0.33900771$ from the ring map (energy independent, as expected
for scaled magnets).

\begin{figure}[htbp]\centering\includegraphics[width=\linewidth]{C:/GIT_REPOS/BELL/Julia/ring/docs/src/assets/acceleration.png}\caption{Acceleration overview (05\_rf\_acceleration.jl). Top left: reference kinetic energy vs turn; dashed line: start of the flat top. Top right: longitudinal phase space at the end of the ramp (turn 2608); black: analytic separatrix H = H(phi\_u,0) of the smooth Hamiltonian with eta0, eta1 at that energy and phi\_s = 150 deg; red cross: unstable fixed point; green: stable fixed point; blue dots: the 60 ensemble particles; orange: test particle \#4 (delta0 = 1e-3) over the last 300 turns (small, near the centre). Bottom left: betatron amplitudes of the x and y test particles vs turn and the prediction A0 sqrt(p\_inj/p0(n)). Bottom right: synchrotron tune; lines: smooth and one-kick formulas evaluated at every turn (step at turn 2608 where phi\_s changes to 180 deg); markers: tracking, tune fitted to delta of particle \#4 in 200-turn windows moved by 50 turns, plotted at the window centre (windows containing the ramp end are skipped).}\end{figure}

\textbf{Animations} (HTML only; the PDF shows the static figure):
\texttt{assets/longitudinal\_ramp.gif} - ensemble every 20 turns during the ramp
and flat top, with the separatrix of the current energy and $\phi_s$; right
panel with $\delta$ divided by the current bucket half-height (removes the
adiabatic shrinking and shows the filamentation).
\texttt{assets/longitudinal\_elements.gif} - one turn at injection, frame = exit of
each element; right panel: change since the turn start (only the cavity
changes $\delta$; $\phi$ changes in every element through $\eta$-like
path/velocity terms).

\begin{quote}\emph{Animation assets/longitudinal\_ramp.gif (see the HTML manual or the GIF file): Animation: ensemble during the ramp (GIF)}\end{quote}

\begin{quote}\emph{Animation assets/longitudinal\_elements.gif (see the HTML manual or the GIF file): Animation: one turn element by element (GIF)}\end{quote}

\clearpage
\section{Symplecticity checks and Taylor-map comparison}

\emph{\texttt{src/Symplectic.jl}, \texttt{src/TaylorMap.jl}, \texttt{src/LieOperators.jl}, script
\texttt{02\_lie\_maps.jl}.}

A map $z\mapsto\mathcal M(z)$ is symplectic iff its Jacobian satisfies
$J^TSJ=S$ with $S=\mathrm{diag}(S_2,S_2,S_2)$, $S_2=\begin{pmatrix}0&1\\-1&0\end{pmatrix}$.
The code reports $\max|J^TSJ-S|$ for

\begin{itemize}
\item every element Lie map (at $z=0$ and at a finite amplitude),
\item the ring map built from Lie factors, the factor-by-factor tracking map
\end{itemize}
  (Jacobian computed exactly by running the tracker on \texttt{TaylorN}), and the
  exact integrator (including the cavity).

Truncated polynomial maps are symplectic only up to the truncation order (the
error grows like $|z|^{\rm MAP\_ORDER}$); the Lie-factor tracker and the exact
integrator are symplectic to round-off.

\textbf{Taylor-map comparison module.} \texttt{TaylorMap.jl} builds the one-turn map
completely independently of the Lie machinery: Hamilton's equations
$dz/ds=S\nabla H$ are integrated with RK4 in the truncated polynomial
algebra (a COSY-style differential-algebra map), element by element, and the
maps are composed. Its coefficients are compared with the Lie-factor product
and with one-turn tracking. Numbers: printed by \texttt{02\_lie\_maps.jl} (see
Sextupoles and the optional optimiser for the tables reproduced from the
latest run).

Latest run of \texttt{02\_lie\_maps.jl} / \texttt{03\_second\_order.jl}:

\begin{center}\small\begin{tabular}{p{0.475\linewidth}p{0.475\linewidth}}\hline
\textbf{check} & \textbf{value}\\\hline
linear ring matrix: Lie product vs product of $\exp(LSA)$ & $2.3\times10^{-13}$\\
linear ring matrix: Taylor/RK4 vs $\exp(LSA)$ & $6.9\times10^{-9}$\\
max 2nd-order coefficient difference Lie vs Taylor (largest coeff 167.5) & $5.1\times10^{-8}$\\
symplectic error ring Taylor map at 0 / at finite amplitude & $6.8\times10^{-13}$ / $1.4\times10^{-5}$\\
one turn: Lie tracking vs Taylor map & $4.7\times10^{-10}$\\
symplectic error of 1-turn Lie-tracking Jacobian & $2.4\times10^{-14}$\\
one turn: exact integrator vs Lie (order 3) / vs Taylor / vs linear & $2.4\times10^{-8}$ / $2.4\times10^{-8}$ / $6.9\times10^{-5}$\\
symplectic error of exact-integrator one-turn Jacobian & $1.1\times10^{-13}$\\
1000 turns Lie tracking: relative spread of $J_x$, $J_y$ & $2.4\times10^{-3}$, $2.7\times10^{-3}$\\
\hline\end{tabular}\end{center}

\clearpage
\section{Sextupoles and the optional optimiser}

\emph{Placement (lattice editing, no physics): \texttt{src/Elements.jl}
(\texttt{find\_optics\_maxima}, \texttt{place\_sextupole\_families}, \texttt{set\_family\_strengths}).
Optimisation: \texttt{src/Optimizers.jl}, called \textbf{only} by
\texttt{scripts/04\_optional\_sextupole\_optimization.jl}.}

**Default: all sextupoles have the strength given in the lattice file, which is
zero for \texttt{magnetic\_2.opt}. No script except 04 changes them.**

\subsection{Placement}

Three families of 0.12 m sextupoles are placed into the nearest existing
sextupole slot or short drift next to local maxima ($\ge0.8$ of the global
maximum) of $D_x$ (family SXD, 2 magnets), $\beta_x$ (SXB, 4) and $\beta_y$
(SYB, 8). A sextupole acts chromatically through $x=D\delta$: the
first-order chromaticity contribution is
$\Delta\xi_{x,y}=\pm\frac{1}{4\pi}\sum k_2L\,\beta_{x,y}D$, so SXD/SXB act mostly
on $\xi_x$ and SYB on $\xi_y$; $D_2$ and $\alpha_1$ are driven by $k_2LD^2$.
The circumference is unchanged.

\subsection{Response (per unit $k_2=1$ m$^{-3}$ for the whole family)}

\begin{center}\small\begin{tabular}{p{0.19\linewidth}p{0.19\linewidth}p{0.19\linewidth}p{0.19\linewidth}p{0.19\linewidth}}\hline
\textbf{family} & \textbf{$\Delta\xi_{1x}$} & \textbf{$\Delta\xi_{1y}$} & \textbf{$\Delta\alpha_1$} & \textbf{$\Delta D_2(0)$ [m]}\\\hline
SXD & 4.32833 & -0.90618 & -1.63792 & 89.86\\
SXB & 3.67887 & -0.74301 & -0.40970 & -0.034\\
SYB & 0.84330 & -4.33701 & -0.24240 & 9.455\\
\hline\end{tabular}\end{center}

$\Delta\eta_1=\Delta\alpha_1$ (sextupoles do not change $c_1,c_2$).

\subsection{Optional example solutions}

These targets are \emph{examples}, not requirements: \textbf{A}
$\xi_{1x}=\xi_{1y}=\alpha_1=0$; \textbf{B} $\xi_{1x}=\xi_{1y}=0$ and minimum rms
$D_2$ (least squares).

\begin{center}\small\begin{tabular}{p{0.238\linewidth}p{0.238\linewidth}p{0.238\linewidth}p{0.238\linewidth}}\hline
\textbf{quantity} & \textbf{natural} & \textbf{A} & \textbf{B}\\\hline
$k_2$ SXD / SXB / SYB [m$^{-3}$] & 0 / 0 / 0 & 0.2304 / 0.2653 / -0.6067 & 0.1456 / 0.3650 / -0.6061\\
$\xi_{1x}$ & -1.4615 & 0.0000 & 0.0000\\
$\xi_{1y}$ & -2.2255 & 0.0000 & 0.0000\\
$\xi_{2x}$ & 8.555 & 3.934 & 5.165\\
$\xi_{2y}$ & 8.104 & -0.717 & 1.236\\
$D_2(0)$ [m] & -10.055 & 4.905 & -2.718\\
max $\vert D_2\vert $ [m] & 14.41 & 10.04 & 3.81\\
rms $D_2$ [m] & 7.558 & 5.728 & 2.411\\
$\alpha_1$ & 0.33901 & 0.00000 & 0.09799\\
$\eta_1$ (270 MeV) & 1.01721 & 0.67820 & 0.77619\\
\hline\end{tabular}\end{center}

Tracking check of A: $\xi_{1x}=0.00001$, $\xi_{2x}=3.934$, $\xi_{1y}=0.00000$,
$\xi_{2y}=-0.717$ (map: identical to 3 digits).

\begin{figure}[htbp]\centering\includegraphics[width=\linewidth]{C:/GIT_REPOS/BELL/Julia/ring/docs/src/assets/chromaticity_corrected.png}\caption{Tunes vs delta after the optional correction A: map (lines) and exact tracking (markers).}\end{figure}

\clearpage
\section{Assumptions and hypotheses}

Every choice that is not a consequence of the lattice file or of the
equations is listed here.

\begin{enumerate}
\item \textbf{Injection energy 10 MeV} (5 MeV/u): not in \texttt{magnetic\_2.opt}; chosen as a typical linac energy and confirmed by the user. Final 270 MeV (135 MeV/u) is from the file header.
\item \textbf{RF parameters} (chosen, not from a design): cavity length $L=1$ m, voltage $V=0.2$ MV, harmonic $h=5$, $\phi_s=150^\circ$ during the ramp and $180^\circ$ on the flat top, one cavity in the first \texttt{OBM} drift. $V$ is large so that the ramp takes \~{}2600 turns instead of $10^5$-$10^6$.
\item \textbf{Cavity field} purely longitudinal on axis, sinusoidal in time, uniform over $L$; \textbf{no transverse RF fields}, no beam loading.
\item **Magnets ramp with $p_0$** (normalised strengths fixed); RF frequency locked to $h\omega_0(p_0)$ of the synchronous particle.
\item \textbf{No fringe fields}, hard-edge elements, sector bends; \textbf{no space charge}, no impedances, no radiation, no errors/misalignments.
\item \textbf{Test amplitudes}: $x_0=2$ mm, $y_0=2$ mm, $\delta_0=10^{-3}$ ($x'=y'=0$); ensemble $\phi\sim N(0,0.6)$, $\delta\sim N(0,0.25\hat\delta)$, 60 particles, seed 1.
\item \textbf{Optimisation targets A and B} are optional examples; the default lattice has sextupoles at zero.
\item \textbf{"Number of periods = 2"} in \texttt{magnetic\_2.opt}: the listed sequence already bends by 360 deg, so it is used \textbf{once} (as in the COSY deck).
\item **$D(0)=0$ in OptiM** initial conditions vs our closed $D(0)=0.0111$ m: OptiM's value is an input, not the closed solution; all other initial Twiss agree to $10^{-5}$.
\item Mass $2\times938.272$ MeV (as OptiM), not the deuteron mass 1875.613 MeV.
\item $\alpha_1$ used for the separatrix is linearly interpolated between the maps at 10 and 270 MeV (it turned out equal at both).
\item Map orders: Hamiltonians to order 4, maps to order 3; tracking with the exact integrator uses $ds=0.2$ m sub-steps (4th order).
\end{enumerate}

\clearpage
\section{Code map}

\begin{center}\small\begin{tabular}{p{0.475\linewidth}p{0.475\linewidth}}\hline
\textbf{file} & \textbf{physics implemented}\\\hline
\texttt{src/Constants.jl} & constants, deuteron, reference particle, $v(\delta)$, $e(\delta)$, exact user <-> canonical transform\\
\texttt{src/Jets.jl} & truncated polynomial space (TaylorSeries \texttt{TaylorN}), orders\\
\texttt{src/Elements.jl} & \textbf{all} element types (Drift, SBend, Quadrupole, Sextupole, thin RFCavity, finite Cavity, Marker), Lattice, lattice-editing helpers (sextupole placement, cavity insertion)\\
\texttt{src/Hamiltonians.jl} & element Hamiltonians as polynomials\\
\texttt{src/LieOperators.jl} & Poisson bracket, $\exp(:f:)$, Dragt–Finn factors $R,f_3,f_4$\\
\texttt{src/OptiMParser.jl} & reader of OptiM \texttt{.opt} files\\
\texttt{src/LinearOptics.jl} & linear maps, closed Twiss, dispersion, tunes, $\alpha_c$\\
\texttt{src/Symplectic.jl} & $J^TSJ-S$ checks\\
\texttt{src/Tracking.jl} & factor-by-factor Lie tracking, thin RF kick\\
\texttt{src/TaylorMap.jl} & independent DA (RK4) one-turn Taylor map, closed orbit, tunes vs $\delta$\\
\texttt{src/NormalForm.jl} & linear normalisation (nonlinear normal form: planned)\\
\texttt{src/SecondOrder.jl} & $D_2$, $\alpha_{0,1}$, $\eta_{0,1}$, $\xi_{1,2}$, longitudinal Hamiltonian, separatrix, $Q_s$\\
\texttt{src/ExactIntegrator.jl} & exact-H symplectic integrator, finite RF cavity, reference renormalisation\\
\texttt{src/Acceleration.jl} & ramp loop, tune fit, element-by-element snapshots\\
\texttt{src/Plotting.jl} & all figures and GIFs\\
\texttt{src/Optimizers.jl} & OPTIONAL sextupole response and example solutions\\
\texttt{scripts/01\_linear\_optics.jl} & linear optics + OptiM/COSY comparison\\
\texttt{scripts/02\_lie\_maps.jl} & Lie maps, symplecticity, Lie vs Taylor, tracking\\
\texttt{scripts/03\_second\_order.jl} & $D_2$, $\alpha/\eta$, chromaticities\\
\texttt{scripts/04\_optional\_sextupole\_optimization.jl} & optional optimiser\\
\texttt{scripts/05\_rf\_acceleration.jl} & RF cavity, ramp, separatrix, GIFs\\
\texttt{scripts/06\_coordinates.jl} & coordinate transformation checks\\
\hline\end{tabular}\end{center}

\clearpage
\section{Programmer manual (stub)}

To be written. Entry point: \texttt{src/LieRing.jl} includes all modules in
dependency order and re-exports their names; scripts do
\texttt{include("../src/LieRing.jl"); using .LieRing}. Build these docs with
\texttt{julia --project=docs docs/make.jl} (HTML in \texttt{docs/build/}) and the PDF with
\texttt{julia docs/make\_pdf.jl} (needs xelatex).

\end{document}
